% Options for packages loaded elsewhere
\PassOptionsToPackage{unicode}{hyperref}
\PassOptionsToPackage{hyphens}{url}
\documentclass[
  ignorenonframetext,
]{beamer}
\newif\ifbibliography
\usepackage{pgfpages}
\setbeamertemplate{caption}[numbered]
\setbeamertemplate{caption label separator}{: }
\setbeamercolor{caption name}{fg=normal text.fg}
\beamertemplatenavigationsymbolsempty
% remove section numbering
\setbeamertemplate{part page}{
  \centering
  \begin{beamercolorbox}[sep=16pt,center]{part title}
    \usebeamerfont{part title}\insertpart\par
  \end{beamercolorbox}
}
\setbeamertemplate{section page}{
  \centering
  \begin{beamercolorbox}[sep=12pt,center]{section title}
    \usebeamerfont{section title}\insertsection\par
  \end{beamercolorbox}
}
\setbeamertemplate{subsection page}{
  \centering
  \begin{beamercolorbox}[sep=8pt,center]{subsection title}
    \usebeamerfont{subsection title}\insertsubsection\par
  \end{beamercolorbox}
}
% Prevent slide breaks in the middle of a paragraph
\widowpenalties 1 10000
\raggedbottom
\AtBeginPart{
  \frame{\partpage}
}
\AtBeginSection{
  \ifbibliography
  \else
    \frame{\sectionpage}
  \fi
}
\AtBeginSubsection{
  \frame{\subsectionpage}
}
\usepackage{iftex}
\ifPDFTeX
  \usepackage[T1]{fontenc}
  \usepackage[utf8]{inputenc}
  \usepackage{textcomp} % provide euro and other symbols
\else % if luatex or xetex
  \usepackage{unicode-math} % this also loads fontspec
  \defaultfontfeatures{Scale=MatchLowercase}
  \defaultfontfeatures[\rmfamily]{Ligatures=TeX,Scale=1}
\fi
\usepackage{lmodern}
\ifPDFTeX\else
  % xetex/luatex font selection
\fi
% Use upquote if available, for straight quotes in verbatim environments
\IfFileExists{upquote.sty}{\usepackage{upquote}}{}
\IfFileExists{microtype.sty}{% use microtype if available
  \usepackage[]{microtype}
  \UseMicrotypeSet[protrusion]{basicmath} % disable protrusion for tt fonts
}{}
\makeatletter
\@ifundefined{KOMAClassName}{% if non-KOMA class
  \IfFileExists{parskip.sty}{%
    \usepackage{parskip}
  }{% else
    \setlength{\parindent}{0pt}
    \setlength{\parskip}{6pt plus 2pt minus 1pt}}
}{% if KOMA class
  \KOMAoptions{parskip=half}}
\makeatother
\usepackage{color}
\usepackage{fancyvrb}
\newcommand{\VerbBar}{|}
\newcommand{\VERB}{\Verb[commandchars=\\\{\}]}
\DefineVerbatimEnvironment{Highlighting}{Verbatim}{commandchars=\\\{\}}
% Add ',fontsize=\small' for more characters per line
\newenvironment{Shaded}{}{}
\newcommand{\AlertTok}[1]{\textcolor[rgb]{1.00,0.00,0.00}{\textbf{#1}}}
\newcommand{\AnnotationTok}[1]{\textcolor[rgb]{0.38,0.63,0.69}{\textbf{\textit{#1}}}}
\newcommand{\AttributeTok}[1]{\textcolor[rgb]{0.49,0.56,0.16}{#1}}
\newcommand{\BaseNTok}[1]{\textcolor[rgb]{0.25,0.63,0.44}{#1}}
\newcommand{\BuiltInTok}[1]{\textcolor[rgb]{0.00,0.50,0.00}{#1}}
\newcommand{\CharTok}[1]{\textcolor[rgb]{0.25,0.44,0.63}{#1}}
\newcommand{\CommentTok}[1]{\textcolor[rgb]{0.38,0.63,0.69}{\textit{#1}}}
\newcommand{\CommentVarTok}[1]{\textcolor[rgb]{0.38,0.63,0.69}{\textbf{\textit{#1}}}}
\newcommand{\ConstantTok}[1]{\textcolor[rgb]{0.53,0.00,0.00}{#1}}
\newcommand{\ControlFlowTok}[1]{\textcolor[rgb]{0.00,0.44,0.13}{\textbf{#1}}}
\newcommand{\DataTypeTok}[1]{\textcolor[rgb]{0.56,0.13,0.00}{#1}}
\newcommand{\DecValTok}[1]{\textcolor[rgb]{0.25,0.63,0.44}{#1}}
\newcommand{\DocumentationTok}[1]{\textcolor[rgb]{0.73,0.13,0.13}{\textit{#1}}}
\newcommand{\ErrorTok}[1]{\textcolor[rgb]{1.00,0.00,0.00}{\textbf{#1}}}
\newcommand{\ExtensionTok}[1]{#1}
\newcommand{\FloatTok}[1]{\textcolor[rgb]{0.25,0.63,0.44}{#1}}
\newcommand{\FunctionTok}[1]{\textcolor[rgb]{0.02,0.16,0.49}{#1}}
\newcommand{\ImportTok}[1]{\textcolor[rgb]{0.00,0.50,0.00}{\textbf{#1}}}
\newcommand{\InformationTok}[1]{\textcolor[rgb]{0.38,0.63,0.69}{\textbf{\textit{#1}}}}
\newcommand{\KeywordTok}[1]{\textcolor[rgb]{0.00,0.44,0.13}{\textbf{#1}}}
\newcommand{\NormalTok}[1]{#1}
\newcommand{\OperatorTok}[1]{\textcolor[rgb]{0.40,0.40,0.40}{#1}}
\newcommand{\OtherTok}[1]{\textcolor[rgb]{0.00,0.44,0.13}{#1}}
\newcommand{\PreprocessorTok}[1]{\textcolor[rgb]{0.74,0.48,0.00}{#1}}
\newcommand{\RegionMarkerTok}[1]{#1}
\newcommand{\SpecialCharTok}[1]{\textcolor[rgb]{0.25,0.44,0.63}{#1}}
\newcommand{\SpecialStringTok}[1]{\textcolor[rgb]{0.73,0.40,0.53}{#1}}
\newcommand{\StringTok}[1]{\textcolor[rgb]{0.25,0.44,0.63}{#1}}
\newcommand{\VariableTok}[1]{\textcolor[rgb]{0.10,0.09,0.49}{#1}}
\newcommand{\VerbatimStringTok}[1]{\textcolor[rgb]{0.25,0.44,0.63}{#1}}
\newcommand{\WarningTok}[1]{\textcolor[rgb]{0.38,0.63,0.69}{\textbf{\textit{#1}}}}
\usepackage{longtable,booktabs,array}
\newcounter{none} % for unnumbered tables
\usepackage{calc} % for calculating minipage widths
\usepackage{caption}
% Make caption package work with longtable
\makeatletter
\def\fnum@table{\tablename~\thetable}
\makeatother
\setlength{\emergencystretch}{3em} % prevent overfull lines
\providecommand{\tightlist}{%
  \setlength{\itemsep}{0pt}\setlength{\parskip}{0pt}}
\usepackage{bookmark}
\IfFileExists{xurl.sty}{\usepackage{xurl}}{} % add URL line breaks if available
\urlstyle{same}
\hypersetup{
  hidelinks,
  pdfcreator={LaTeX via pandoc}}

\author{\texorpdfstring{}{}}
\date{}

\begin{document}

\section{Reproducibility}\label{reproducibility}

\begin{frame}[fragile]{Software}
\protect\phantomsection\label{software}
\begin{itemize}
\tightlist
\item
  Python \(\geq 3.10\), NumPy, Matplotlib, PyYAML (declared in
  \texttt{projects/special\_number\_proximity/pyproject.toml}; root
  \texttt{uv\ sync} supplies the environment).
\item
  Tests:
  \texttt{uv\ run\ pytest\ projects/special\_number\_proximity/tests/\ -\/-cov=projects/special\_number\_proximity/src\ -\/-cov-fail-under=90}.
\end{itemize}
\end{frame}

\begin{frame}[fragile]{Regenerating artifacts}
\protect\phantomsection\label{regenerating-artifacts}
\begin{Shaded}
\begin{Highlighting}[]
\ExtensionTok{uv}\NormalTok{ run python projects/special\_number\_proximity/scripts/proximity\_monte\_carlo.py}
\CommentTok{\# Optional: alternate project root (e.g. temporary checkout)}
\ExtensionTok{uv}\NormalTok{ run python projects/special\_number\_proximity/scripts/proximity\_monte\_carlo.py }\AttributeTok{{-}{-}project{-}dir}\NormalTok{ /path/to/projects/special\_number\_proximity}
\end{Highlighting}
\end{Shaded}

The script prints absolute paths to \textbf{five} outputs: JSON summary,
CSV constant table, and three PNG figures
(\texttt{proximity\_histogram.png},
\texttt{proximity\_histogram\_pooled.png},
\texttt{proximity\_histogram\_mu.png}).

Edits to \texttt{manuscript/config.yaml} under \texttt{experiment:}
change seeds, sample sizes, \texttt{q\_max},
\texttt{quadratic\_candidates}, \texttt{compare\_mod1},
\texttt{q\_sensitivity}, \texttt{histogram\_uniform\_cap}, and Beta
parameters without editing Python.
\end{frame}

\begin{frame}[fragile]{Data paths}
\protect\phantomsection\label{data-paths}
{\def\LTcaptype{none} % do not increment counter
\begin{longtable}[]{@{}
  >{\raggedright\arraybackslash}p{(\linewidth - 2\tabcolsep) * \real{0.6250}}
  >{\raggedright\arraybackslash}p{(\linewidth - 2\tabcolsep) * \real{0.3750}}@{}}
\toprule\noalign{}
\begin{minipage}[b]{\linewidth}\raggedright
Artifact
\end{minipage} & \begin{minipage}[b]{\linewidth}\raggedright
Path
\end{minipage} \\
\midrule\noalign{}
\endhead
Summary JSON &
\texttt{projects/special\_number\_proximity/output/data/proximity\_summary.json} \\
Table CSV &
\texttt{projects/special\_number\_proximity/output/data/proximity\_constants.csv} \\
Uniform \(\delta_Q\) histogram &
\texttt{projects/special\_number\_proximity/output/figures/proximity\_histogram.png} \\
Pooled \(\delta_Q\) histogram &
\texttt{projects/special\_number\_proximity/output/figures/proximity\_histogram\_pooled.png} \\
Pooled \(\mu_Q\) histogram &
\texttt{projects/special\_number\_proximity/output/figures/proximity\_histogram\_mu.png} \\
Lattice cross-check &
\texttt{projects/special\_number\_proximity/output/data/lattice\_crosscheck.json} \\
\bottomrule\noalign{}
\end{longtable}
}

Run
\texttt{uv\ run\ python\ projects/special\_number\_proximity/scripts/02\_lattice\_crosscheck.py}
to refresh the lattice JSON.
\end{frame}

\begin{frame}[fragile]{Combined manuscript build}
\protect\phantomsection\label{combined-manuscript-build}
\begin{Shaded}
\begin{Highlighting}[]
\ExtensionTok{uv}\NormalTok{ run python scripts/03\_render\_pdf.py }\AttributeTok{{-}{-}project}\NormalTok{ special\_number\_proximity}
\ExtensionTok{python3} \AttributeTok{{-}m}\NormalTok{ infrastructure.validation.cli markdown projects/special\_number\_proximity/manuscript/}
\end{Highlighting}
\end{Shaded}

Authoring conventions for Markdown and cross-refs:
\href{../docs/manuscript_conventions.md}{\texttt{docs/manuscript\_conventions.md}}.
\end{frame}

\end{document}
